\documentclass[tikz,border=6pt]{standalone}
\usepackage[no-math]{fontspec}
\usepackage[brazil]{babel}
\usepackage{amsmath,amssymb}
\IfFileExists{C:/Windows/Fonts/Lato_Regular.ttf}{%
  \setmainfont{Lato}[
    Path = C:/Windows/Fonts/,
    UprightFont = Lato_Regular.ttf,
    BoldFont = Lato_Bold.ttf,
    ItalicFont = Lato_Italic.ttf,
    BoldItalicFont = Lato_BoldItalic.ttf
  ]
  \setsansfont{Lato}[
    Path = C:/Windows/Fonts/,
    UprightFont = Lato_Regular.ttf,
    BoldFont = Lato_Bold.ttf,
    ItalicFont = Lato_Italic.ttf,
    BoldItalicFont = Lato_BoldItalic.ttf
  ]
}{%
  \RequirePackage[default]{lato}
}
\renewcommand{\familydefault}{\sfdefault}

\usepackage{xcolor}
\usepackage{tikz}
\usepackage{pgfplots}
\pgfplotsset{compat=1.18}
\usetikzlibrary{arrows.meta,positioning,shapes.geometric,fit,backgrounds,calc,decorations.pathreplacing,patterns}

% Paleta institucional UEFS / CHF662
\definecolor{uefsNavy}{HTML}{1B2A4A}
\definecolor{uefsTeal}{HTML}{0E7C7B}
\definecolor{uefsForest}{HTML}{2D6A4F}
\definecolor{uefsGold}{HTML}{C4983B}
\definecolor{uefsGoldDark}{HTML}{8F6B1E}
\definecolor{uefsAlert}{HTML}{B02A37}
\definecolor{uefsOrange}{HTML}{D95F02}
\definecolor{uefsInk}{HTML}{22252A}
\definecolor{uefsMuted}{HTML}{5C6470}
\definecolor{uefsWater}{HTML}{2E86AB}
\definecolor{uefsSand}{HTML}{D9B382}
\definecolor{bgNavy}{HTML}{F2F5F9}
\definecolor{bgTeal}{HTML}{E9F5F4}
\definecolor{bgForest}{HTML}{EEF5F0}
\definecolor{bgGold}{HTML}{FCF7EB}
\definecolor{bgAlert}{HTML}{FDF2F2}
\definecolor{bgGray}{HTML}{F8FAFC}
\definecolor{lineGray}{HTML}{CFD8DC}

\tikzset{
  >=Stealth,
  aaCard/.style={draw=#1, fill=white, rounded corners=4pt, line width=1pt, inner sep=6pt, align=center, font=\small},
  aaBox/.style={draw=#1, rounded corners=4pt, line width=1.1pt, inner sep=6pt, align=center, font=\footnotesize\sffamily, text=uefsInk},
  aaTitle/.style={font=\normalsize\bfseries\sffamily, text=uefsNavy, align=center},
  aaSub/.style={font=\footnotesize\sffamily, text=uefsMuted, align=center},
  aaArrow/.style={->, >=Stealth, line width=1.4pt},
  aaDArrow/.style={<->, >=Stealth, line width=1.4pt}
}

\usepgfplotslibrary{fillbetween}

\begin{document}
\begin{tikzpicture}
\begin{axis}[
  width=15.6cm,
  height=7.8cm,
  title={Morfopedologia: Relação Relevo, Solo e Água ao Longo de uma Vertente},
  title style={font=\normalsize\bfseries\sffamily, text=uefsNavy, yshift=4pt},
  axis lines=left,
  axis line style={lineGray, line width=1pt},
  tick style={uefsMuted},
  xmin=0, xmax=10.2,
  ymin=0, ymax=6.7,
  xtick={0, 2, 4, 6, 8, 10},
  ytick={0, 1, 2, 3, 4, 5, 6},
  tick label style={font=\footnotesize\sffamily, text=uefsInk},
  xlabel={\textbf{Distância ao longo da vertente (topossequência)}},
  ylabel={\textbf{Altitude relativa (m)}},
  label style={font=\footnotesize\sffamily, text=uefsInk},
  legend style={
    at={(0.02,0.04)}, anchor=south west,
    font=\scriptsize\sffamily, fill=white, fill opacity=0.92,
    draw=lineGray, rounded corners=3pt, inner sep=4pt,
    cells={anchor=west}
  },
  legend image post style={scale=0.8},
  clip=true
]

% Superfície topográfica côncavo-convexa e base do manto de solo
\addplot[name path=base, draw=none, domain=0:10] {0};

\addplot[
  name path=bed,
  dashed, color=uefsNavy, line width=1.1pt,
  domain=0:10, samples=120, forget plot
] {
  (3.6 - 0.25*x - 1.2/(1 + exp(-(x - 4.3)*1.1)) + 0.75)
  - (0.65 - 0.32*exp(-((x - 4.5)^2)/2.2) + 0.35*exp(-((x - 9.0)^2)/1.5))
};

\addplot[
  name path=surf,
  color=uefsForest, line width=2.2pt,
  domain=0:10, samples=120, forget plot
] {3.6 - 0.25*x - 1.2/(1 + exp(-(x - 4.3)*1.1)) + 0.75};

% Preenchimento do substrato (Rocha / Saprólito)
\addplot[fill=lineGray!55, draw=none] fill between[of=base and bed];
\addlegendentry{Rocha / Saprólito (substrato)}

% Preenchimento do manto de solo
\addplot[fill=uefsSand!78, draw=none] fill between[of=bed and surf];
\addlegendentry{Manto de solo (varia em espessura na vertente)}

% Re-traçar linhas para nitidez sobre os preenchimentos
\addplot[dashed, color=uefsNavy, line width=1.1pt, domain=0:10, samples=120, forget plot] {
  (3.6 - 0.25*x - 1.2/(1 + exp(-(x - 4.3)*1.1)) + 0.75)
  - (0.65 - 0.32*exp(-((x - 4.5)^2)/2.2) + 0.35*exp(-((x - 9.0)^2)/1.5))
};
\addplot[color=uefsForest, line width=2.3pt, domain=0:10, samples=120, forget plot]
  {3.6 - 0.25*x - 1.2/(1 + exp(-(x - 4.3)*1.1)) + 0.75};

% Canal fluvial no fundo de vale
\fill[uefsWater!80] (axis cs:9.05, 0.56) rectangle (axis cs:10.0, 0.90);
\draw[uefsWater, line width=1pt] (axis cs:9.05, 0.90) -- (axis cs:10.0, 0.90);
\node[font=\scriptsize\bfseries\sffamily, text=uefsWater] at (axis cs:9.52, 1.12) {Rio};

% Seta de infiltração no topo
\draw[aaArrow, color=uefsWater, line width=1.5pt] (axis cs:1.2, 3.85) -- (axis cs:1.2, 3.12);
\node[font=\scriptsize\bfseries\sffamily, text=uefsWater, anchor=west]
  at (axis cs:1.32, 3.45) {Infiltração};

% Seta de escoamento superficial na média vertente
\draw[aaArrow, color=uefsAlert, line width=1.8pt] (axis cs:3.6, 3.22) -- (axis cs:5.25, 2.28);
\node[font=\scriptsize\bfseries\sffamily, text=uefsAlert, align=center, rotate=-25]
  at (axis cs:4.72, 2.98) {Escoamento superficial};

% 4 caixas descritivas posicionadas CLARAMENTE ACIMA da superfície com linhas-guia sutis
\node[aaBox=uefsForest, fill=white, fill opacity=0.95, text width=2.55cm, inner sep=4pt, font=\scriptsize\sffamily]
  (z1) at (axis cs:1.25, 5.55)
  {\textbf{\color{uefsForest}Topo (Interflúvio)}\\[2pt]
   Declividade baixa\\Infiltração alta\\Solos profundos\\(Latossolos)};

\node[aaBox=uefsAlert, fill=white, fill opacity=0.95, text width=2.95cm, inner sep=4pt, font=\scriptsize\sffamily]
  (z2) at (axis cs:3.95, 4.95)
  {\textbf{\color{uefsAlert}Alta / Média Vertente}\\[2pt]
   Declividade acentuada\\Escoamento veloz\\Solos rasos / pedregosos\\(Neossolos / Cambissolos)};

\node[aaBox=uefsGoldDark, fill=white, fill opacity=0.95, text width=2.8cm, inner sep=4pt, font=\scriptsize\sffamily]
  (z3) at (axis cs:6.65, 3.85)
  {\textbf{\color{uefsGoldDark}Baixa Vertente}\\[2pt]
   Sopé côncavo\\Acúmulo de água e colúvios\\(Argissolos / Planossolos)};

\node[aaBox=uefsNavy, fill=white, fill opacity=0.95, text width=2.75cm, inner sep=4pt, font=\scriptsize\sffamily]
  (z4) at (axis cs:8.95, 2.70)
  {\textbf{\color{uefsNavy}Fundo de Vale / Rio}\\[2pt]
   Planície fluvial\\Solos hidromórficos\\(Gleissolos)\\Risco de assoreamento};

% Linhas-guia sutis até os respectivos compartimentos da vertente
\draw[densely dotted, color=uefsForest, line width=0.9pt] (z1.south) -- (axis cs:1.25, 4.02);
\draw[densely dotted, color=uefsAlert, line width=0.9pt] (z2.south) -- (axis cs:3.95, 2.95);
\draw[densely dotted, color=uefsGoldDark, line width=0.9pt] (z3.south) -- (axis cs:6.65, 1.58);
\draw[densely dotted, color=uefsNavy, line width=0.9pt] (z4.south) -- (axis cs:9.10, 0.92);

\end{axis}
\end{tikzpicture}
\end{document}
